% Prinzip des Doppelsteppstichs nach der bisherigen gemeinfreien Grafik.
\begin{tikzpicture}[x=1cm,y=1cm,line cap=round,line join=round,
  ober/.style={draw=fablab,line width=1.35pt},
  unter/.style={draw=orange!80!black,line width=1.35pt},
  beschriftung/.style={font=\small,anchor=west}]
\node[beschriftung,font=\small\bfseries] at (0,4.7){Stichbildung (schematisch)};
\path[draw=black!40,fill=black!6] (0,2.55) rectangle (9,3.05);
\path[draw=black!40,fill=black!12] (0,2.05) rectangle (9,2.55);
\node[font=\scriptsize,fill=white,inner sep=2pt] at (8.0,2.8){Stofflagen};
% Bereits gebildete Stiche: Verschlingung zwischen den Stofflagen.
\foreach \x in {.75,1.85,2.95}{
 \draw[ober] (\x-.55,3.05)--(\x-.14,3.05) .. controls (\x-.02,3.05) and (\x-.16,2.47) .. (\x,2.47)
 .. controls (\x+.17,2.47) and (\x+.02,3.05) .. (\x+.16,3.05)--(\x+.55,3.05);
 \draw[unter] (\x-.55,2.05)--(\x-.15,2.05) .. controls (\x,2.05) and (\x-.16,2.62) .. (\x,2.62)
 .. controls (\x+.16,2.62) and (\x,2.05) .. (\x+.15,2.05)--(\x+.55,2.05);
}
% Nadel mit Öhr, Oberfadenschlinge und Greifer darunter.
\path[draw=black!75,fill=white,line width=.8pt] (4.47,4.3)--(4.47,2.05)--(4.6,1.7)--(4.73,2.05)--(4.73,4.3)--cycle;
\draw[black!60] (4.6,1.94) ellipse (.045 and .10);
\draw[ober] (3.5,3.05)--(4.32,3.05) .. controls (4.42,2.65) and (4.25,1.1) .. (4.9,.8)
 .. controls (5.65,.4) and (6.28,1.12) .. (5.97,1.53)
 .. controls (5.7,1.9) and (4.78,1.94) .. (4.65,1.94)--(4.88,4.28);
\path[draw=black!65,fill=black!6,line width=.8pt] (5.85,1.45)
 .. controls (6.5,1.95) and (7.95,1.72) .. (8.5,1.2)
 .. controls (7.9,.7) and (6.48,.72) .. (5.85,1.45);
\draw[-{Stealth},black!65] (6.65,1.2)--(7.65,1.2);
\draw[-{Stealth},black!65] (4.15,3.5)--(4.15,4.18);
\node[beschriftung] at (5.1,4.0){Nadel}; \draw[black!45] (5.0,3.92)--(4.73,3.8);
\node[beschriftung] at (6.35,.25){Greifer}; \draw[black!45] (6.3,.4)--(6.25,1.02);
\node[beschriftung,text=fablab] at (0,.6){Oberfadenschlinge}; \draw[fablab!70] (3.55,.66)--(4.53,.95);
% Fertige Naht mit identifizierbaren Ober- und Unterfäden.
\node[beschriftung,font=\small\bfseries] at (0,-.65){Fertige Naht: Verschlingung im Stoff};
\path[draw=black!40,fill=black!6] (0,-1.75) rectangle (9,-1.25);
\path[draw=black!40,fill=black!12] (0,-2.25) rectangle (9,-1.75);
\foreach \x in {.7,1.8,2.9,4.0,5.1,6.2,7.3,8.4}{
 \draw[ober] (\x-.55,-1.25)--(\x-.14,-1.25) .. controls (\x-.02,-1.25) and (\x-.16,-1.83) .. (\x,-1.83)
 .. controls (\x+.17,-1.83) and (\x+.02,-1.25) .. (\x+.16,-1.25)--(\x+.55,-1.25);
 \draw[unter] (\x-.55,-2.25)--(\x-.15,-2.25) .. controls (\x,-2.25) and (\x-.16,-1.68) .. (\x,-1.68)
 .. controls (\x+.16,-1.68) and (\x,-2.25) .. (\x+.15,-2.25)--(\x+.55,-2.25);
}
\draw[ober] (.2,-2.9)--(.8,-2.9); \node[beschriftung] at (1,-2.9){Oberfaden};
\draw[unter] (4.3,-2.9)--(4.9,-2.9); \node[beschriftung] at (5.1,-2.9){Unterfaden};
\end{tikzpicture}
